\chapter*{Appendix B: Common Surgical Incisions \& Access}
\addcontentsline{toc}{chapter}{Appendix B: Common Surgical Incisions \& Access}
\markboth{Appendix B: Surgical Incisions}{Appendix B: Surgical Incisions}

The choice of surgical incision is determined by the need for adequate operative exposure, the speed of access required, cosmetic preferences, and the avoidance of nerves and major vessels. Proper technique during entry and closure minimizes postoperative pain and the risk of incisional hernia.

\section*{1. Midline Incision (Laparotomy)}
\textbf{Description:} An incision made along the linea alba, extending from the xiphoid process to the pubic symphysis (upper midline, lower midline, or full midline).
\begin{itemize}
  \item \textbf{Anatomy:} Follows the avascular linea alba, avoiding muscle fibers. It passes through the skin, subcutaneous fat, linea alba, transversalis fascia, and peritoneum.
  \item \textbf{Clinical Exposure:} Provides rapid, wide exposure to the entire abdominal cavity. Preferred in trauma, emergency laparotomies, and multi-quadrant surgeries.
  \item \textbf{Safety/Closure:} Closed using a continuous loop absorbable or non-absorbable suture with a 4:1 suture-to-wound length ratio. Higher risk of incisional hernia compared to transverse incisions.
\end{itemize}

\section*{2. Paramedian Incision}
\textbf{Description:} A vertical incision made parallel to the midline (typically 2-4 cm lateral) over the rectus abdominis muscle.
\begin{itemize}
  \item \textbf{Anatomy:} The anterior rectus sheath is incised, the rectus muscle is retracted laterally (or split), and the posterior rectus sheath and peritoneum are incised.
  \item \textbf{Clinical Exposure:} Historically used for upper abdominal exposures (e.g., splenectomy, gastrectomy). Rectus muscle acts as a buttress upon closure.
  \item \textbf{Safety/Closure:} Less prone to hernia than midline incisions, but takes longer to open and close, and carries a risk of injuring the intercostal nerves and epigastric vessels.
\end{itemize}

\section*{3. Subcostal Incision (Kocher's)}
\textbf{Description:} An oblique incision made 2-5 cm below and parallel to the costal margin. A bilateral subcostal incision is known as a Chevron incision.
\begin{itemize}
  \item \textbf{Anatomy:} Cuts through the anterior rectus sheath, rectus abdominis muscle, posterior rectus sheath, internal oblique, transversus abdominis, and peritoneum.
  \item \textbf{Clinical Exposure:} 
  \begin{itemize}
    \item Right subcostal (Kocher's): Open cholecystectomy, biliary tract surgery, and liver resections.
    \item Left subcostal: Open splenectomy.
    \item Chevron (Bilateral): Liver transplantation and pancreaticoduodenectomy (Whipple).
  \end{itemize}
  \item \textbf{Safety/Closure:} Protect the intercostal nerves to prevent postoperative abdominal wall paresis or rectus atrophy. Excellent strength with low hernia rates.
\end{itemize}

\section*{4. McBurney's and Rocky-Davis Incisions}
\textbf{Description:} Classic incisions used in the right lower quadrant for appendectomy.
\begin{itemize}
  \item \textbf{Anatomy:} 
  \begin{itemize}
    \item McBurney's: An oblique, muscle-splitting incision centered at McBurney's point (one-third from ASIS to umbilicus).
    \item Rocky-Davis: A transverse, muscle-splitting incision at the same location, offering better cosmesis.
  \end{itemize}
  \item \textbf{Access:} Vermiform appendix and cecum.
  \item \textbf{Safety/Closure:} Muscles (external oblique, internal oblique, transversus abdominis) are split in the direction of their fibers rather than cut, preserving abdominal wall integrity.
\end{itemize}

\section*{5. Pfannenstiel Incision}
\textbf{Description:} A low, transverse curved incision made 2-5 cm above the pubic symphysis, typically within the pubic hairline.
\begin{itemize}
  \item \textbf{Anatomy:} Skin and subcutaneous tissues are cut transversely. The anterior rectus sheath is incised transversely and dissected off the rectus muscles. The rectus muscles are separated vertically along the midline, and the peritoneum is entered vertically.
  \item \textbf{Clinical Exposure:} Pelvic cavity, uterus, bladder, and adnexa.
  \item \textbf{Common Procedures:} Cesarean section, abdominal hysterectomy, and pelvic reconstruction.
  \item \textbf{Safety/Closure:} Exceptional cosmetic result and high wound strength. Guard against injury to the iliohypogastric and ilioinguinal nerves.
\end{itemize}

\section*{6. Median Sternotomy}
\textbf{Description:} A vertical midline incision over the sternum.
\begin{itemize}
  \item \textbf{Anatomy:} Skin and subcutaneous tissues are divided, and the sternum is split vertically along its midline using a sternal saw.
  \item \textbf{Clinical Exposure:} Mediastinum, heart, and great vessels.
  \item \textbf{Common Procedures:} Coronary artery bypass grafting (CABG) and heart valve replacements.
  \item \textbf{Safety/Closure:} Closed with stainless steel wires looped around or through the sternum. Sternal dehiscence and mediastinitis are rare but severe complications.
\end{itemize}

\clearpage
